\documentclass[10pt,a4paper]{amsart}
\usepackage[margin=19mm]{geometry}
\usepackage{amsmath,amssymb,mathtools}
\usepackage[scr=rsfs,cal=euler]{mathalfa}
\usepackage{xfrac}
\usepackage[unicode,hidelinks,bookmarks=false]{hyperref}
\newcommand{\ie}{i.e.\ }
\newcommand{\cf}{cf.\ }
\newcommand{\resp}{respectively\ }
\newcommand{\Subsection}{\subsection}
\providecommand{\nobreakdash}{\nobreak\hbox{-}}
\setlength{\emergencystretch}{2em}
\multlinegap0pt
\usepackage{amssymb}
\newtheorem{lem}{Lemma}
\title{On the second Bohr radius for vector valued pluriharmonic functions}
\author{Himadri Halder}
\date{Source: arXiv:2601.12608v2 (CC BY 4.0)}
\begin{document}
\maketitle

\noindent\emph{Source and licence.} On the second Bohr radius for vector valued pluriharmonic functions. Version arXiv:2601.12608v2 (CC BY 4.0), distributed by arXiv under the Creative Commons Attribution 4.0 International licence, http://creativecommons.org/licenses/by/4.0/, as recorded in the arXiv licence metadata for this exact version and shown on the abstract page https://arxiv.org/abs/2601.12608v2. Full licence notice: \texttt{licenses/arxiv-2601.12608-CC-BY-4.0.txt}.\par\smallskip

\noindent\textit{Editorial scope.} Counted unit: the article's lem lem-3.5 (source0.tex lines 354--357). The article's complete original proof of it is delivered with it (source0.tex lines 358--443, one \texttt{proof} environment of 86 lines). No mathematics is written, repaired or re-derived: the statement and the proof are the article's own lines, in the article's own notation, and no retained line is substituted or re-flowed.  Retained uncounted as the source-local context the proof needs: lines 272--274.  Nothing was omitted from any retained span, so the package carries no omission record.  No retained line contains a citation, so the wrapper carries no bibliography and no bibliography entry.  No retained line contains a figure environment, an image-inclusion command or drawing code, so no image is delivered and the package declares no asset dependency.  Editorial formatting only, in addition: an AMS \texttt{amsart} preamble that restates the article's own declarations and macros used by the retained lines, namely the environments \texttt{lem} on separate counters, together with the standard packages those lines need.  The retained environments print in this document as Lemma 1 in order of appearance, whereas the article prints the same items with the numbering of its own full document.  The complete native source of the article as distributed in the arXiv e-print for this exact version is bundled unchanged as \texttt{sources/192965/source0.tex}.
	 \begin{equation} \label{e-1.3-a}
	 	f(z)=\sum_{m=0}^{\infty} \sum_{|\alpha|=m} a_{\alpha}\, z^{\alpha} + \sum_{m=1}^{\infty} \sum_{|\alpha|=m} b^{*}_{\alpha}\, \bar{z}^{\alpha},
	 \end{equation}

	 \begin{lem} \label{lem-3.5}
		Let $\Omega_{1}$ and $\Omega_{2}$ be two bounded simply connected complete Reinhardt domains in $\mathbb{C}^n$. Then we have $$\frac{L_{\lambda}(\Omega_{2}, p,U)}{S(\Omega_{1},\Omega_{2})S(\Omega_{2},\Omega_{1})} \, \leq L_{\lambda}(\Omega_{1}, p,U) \leq S(\Omega_{1},\Omega_{2})S(\Omega_{2},\Omega_{1}) \, L_{\lambda}(\Omega_{2}, p,U).
		$$
	\end{lem}

\begin{proof}
We first prove the right hand inequality.
Let $f \in \mathcal{PH}(\Omega_2,\mathcal{B}(\mathcal{H}))$ be of the form \eqref{e-1.3-a} with $\|f\|_{\Omega_2} \le 1$.

Let $\delta_1>0$.  
By the definition of $S(\Omega_1,\Omega_2)$,
\[
z \in \Omega_1
\quad \Rightarrow \quad
\frac{z}{S(\Omega_1,\Omega_2)+\delta_1} \in \Omega_2.
\]

Define
\[
g(z)=f\!\left(
\frac{z}{S(\Omega_1,\Omega_2)+\delta_1}
\right).
\]
Then $g \in \mathcal{PH}(\Omega_1,\mathcal{B}(\mathcal{H}))$
and $\|g\|_{\Omega_1}\le 1$.
\par
Now, for fix $0<\varepsilon<L_{\lambda}(\Omega_1,p,U)$, let $r = L_{\lambda}(\Omega_1,p,U)-\varepsilon$.
Since $r<L_{\lambda}(\Omega_1,p,U)$, we have
\[
\sum_{m=0}^{\infty}\sum_{|\alpha|=m}
\sup_{z\in r\Omega_1}
(\|U(a_\alpha)\|^p+\|U(b_\alpha)\|^p)
\left|
\left(
\frac{z}{S(\Omega_1,\Omega_2)+\delta_1}
\right)^\alpha
\right|^p
\le \lambda^p.
\]

Hence
\begin{equation} \label{e-compare-1}
\sum_{m=0}^{\infty}\sum_{|\alpha|=m}
\sup_{z\in \frac{r}{S(\Omega_1,\Omega_2)+\delta_1}\Omega_1}
(\|U(a_\alpha)\|^p+\|U(b_\alpha)\|^p)
|z^\alpha|^p
\le \lambda^p.
\end{equation}

Now let $\delta_2>0$.  
Using the definition of $S(\Omega_2,\Omega_1)$, we have
\[
z \in \Omega_2
\quad \Rightarrow \quad
\frac{z}{S(\Omega_2,\Omega_1)+\delta_2} \in \Omega_1.
\]
Thus, from \eqref{e-compare-1}, we obtain
\[
\sum_{m=0}^{\infty}\sum_{|\alpha|=m}
\sup_{z\in
	\frac{r}{(S(\Omega_1,\Omega_2)+\delta_1)(S(\Omega_2,\Omega_1)+\delta_2)}
	\Omega_2}
(\|U(a_\alpha)\|^p+\|U(b_\alpha)\|^p)
|z^\alpha|^p
\le \lambda^p.
\]

By definition of $L_{\lambda}(\Omega_2,p,U)$,
\[
\frac{r}
{(S(\Omega_1,\Omega_2)+\delta_1)(S(\Omega_2,\Omega_1)+\delta_2)}
\le
L_{\lambda}(\Omega_2,p,U).
\]
Letting $\delta_1,\delta_2 \to 0$ gives
\[
r
\le
S(\Omega_1,\Omega_2)S(\Omega_2,\Omega_1)
L_{\lambda}(\Omega_2,p,U).
\]

Since $\varepsilon$ is arbitrary, we deduce that
\[
L_{\lambda}(\Omega_1,p,U)
\le
S(\Omega_1,\Omega_2)S(\Omega_2,\Omega_1)
L_{\lambda}(\Omega_2,p,U).
\]
Interchanging $\Omega_1$ and $\Omega_2$ in the preceding inequality gives the left hand inequality. 
\end{proof}

\end{document}
